\chapter{Algoritmos y estructuras de apoyo}

\section{Propósito}

Proveer una base estable y tipada para implementar algoritmos de grafos y construir el grafo de conflictos. El uso de \textbf{TypeScript} facilita genéricos, interfaces y trazabilidad de tipos.

\section{Estructuras fundamentales}

\subsection{\texttt{Vertex\textless T\textgreater}}
\begin{lstlisting}[language=TypeScript,caption={Definición simplificada de Vertex\textless T\textgreater}]
export class Vertex<T> {
  constructor(public readonly id: string, private value: T, private color?: string) {}
  getValue(): T { return this.value; }
  setColor(c: string) { this.color = c; }
  getColor(): string | undefined { return this.color; }
    \subsection{Visión general y contrato base}

    Las restricciones (constraints) del sistema se modelan como clases que implementan contratos abstractos. Para mantener claridad se separan dos ámbitos:

    \begin{itemize}
      \item Allocator (encargado de asignar franjas y días a las materias): restricciones de \textbf{día} y restricciones de \textbf{franja} (slot).
      \item Student (encargado de verificar horarios por estudiante): restricciones aplicadas al horario final del estudiante.
    \end{itemize}

    Los contratos base son los siguientes:
    \begin{lstlisting}[language=TypeScript,caption={Contratos base para restricciones de asignación}]
    export abstract class DayConstraint {
      abstract name: string;
      abstract isSatisfied(day: string, materia: MateriaInput, assigned: FranjaHoraria[], ctx: DayConstraintContext): boolean;
    }

    export abstract class SlotConstraint {
      abstract name: string;
      abstract isSatisfied(franja: FranjaHoraria, materia: MateriaInput, assigned: FranjaHoraria[], strict: boolean): boolean;
    }
    \end{lstlisting}

    Y para restricciones sobre horarios de estudiantes:
    \begin{lstlisting}[language=TypeScript,caption={Contrato base para restricciones de horarios estudiantiles}]
    export abstract class ScheduleConstraint {
      abstract name: string;
      abstract isSatisfied(candidate: AsignacionResultado, confirmed: AsignacionResultado[], grupoElegido: Map<string, string>): boolean;
    }
    \end{lstlisting}

    \subsection{Listado de restricciones (resumen)}

    Para facilitar la lectura listamos las restricciones agrupadas por ámbito y tipo. Las referencias entre corchetes corresponden a los códigos usados en el repositorio.

    \paragraph{Allocator - Day constraints}
    \begin{itemize}
      \item \texttt{NoRepeatedDay [R2]}: Un mismo grupo no puede tener dos bloques el mismo día.
      \item \texttt{NoConsecutiveDay [R3]}: Materias de 3 créditos no pueden tener días consecutivos.
      \item \texttt{MinDayDistance [R5]}: Materias de 4 créditos requieren distancia mínima entre días.
      \item \texttt{NoHeavyDaySaturation [R8]}: Evitar saturar un día con materias pesadas del mismo semestre/carrera.
    \end{itemize}

    \paragraph{Allocator - Slot constraints}
    \begin{itemize}
      \item \texttt{NoRepeatedSlotNumber [R4]}: Materias de 2 créditos no repiten número de franja.
      \item \texttt{LabPrefersAfternoon [R7]}: Laboratorios prefieren franjas de tarde (aplicable en pasada estricta).
      \item \texttt{HeavyHourRange [R6]}: Materias 3+ créditos solo entre horas permitidas (ej. 08:00-20:00).
      \item \texttt{NoConsecutiveSlot [R9]}: No permitir franjas numéricamente consecutivas dentro del mismo día.
    \end{itemize}

    \paragraph{Student constraints}
    \begin{itemize}
      \item \texttt{NoOverlap}: Evitar solapamientos temporales entre bloques confirmados.
      \item \texttt{NoGroupRepetition}: Un mismo grupo no puede aparecer dos veces en el mismo día.
      \item \texttt{ConsistentGroup}: Si se eligió un grupo para una materia, aceptar sólo ese grupo.
      \item \texttt{NoImmediateSedeChange}: No cambiar de sede en bloques inmediatamente consecutivos.
      \item \texttt{MaxSubjectsPerDay}: Límite de materias distintas por día en el horario del estudiante.
      \item \texttt{MaxFatigue}: Límite de "fatiga" (distancia) entre bloques según un mapa de distancias.
    \end{itemize}

\subsection{Diagrama de clases (restricciones)}

\begin{figure}[h!]
  \centering
  \includegraphics[width=0.9\textwidth]{images/constraints-class-diagram.png}
  \caption{Diagrama de clases resumen de las restricciones (abstractas y concretas) — tanto para asignación (Allocator) como para horarios estudiantiles (Student)}
  \label{fig:constraints-class-diagram}
\end{figure}

\subsection{Uso de imágenes (rutas y URLs)}

El informe referencia diagramas e imágenes que pueden servirse desde rutas relativas locales o URLs remotas. En este repositorio utilizamos rutas relativas (referenciando la carpeta \texttt{images/}) para asegurar reproducibilidad; si requieres servirlas desde URLs remotas, basta cambiar las referencias de forma sistemática. Ejemplo de URL alternativa:

\begin{verbatim}
\includegraphics[width=0.9\textwidth]{https://mi-servidor.org/diagrams/constraints-class-diagram.png}
\end{verbatim}

    \subsection{Implementación: verificador (Allocator)}

    El verificador de allocator agrupa reglas de día y de franja y las evalúa antes de aceptar una asignación:
    \begin{lstlisting}[language=TypeScript,caption={Verificador de restricciones (Allocator)}]
    export class AllocatorConstraintChecker {
      constructor(private readonly dayConstraints: DayConstraint[], private readonly slotConstraints: SlotConstraint[]) {}

      checkDay(day: string, materia: MateriaInput, assigned: FranjaHoraria[], ctx: DayConstraintContext) {
        for (const c of this.dayConstraints) {
          if (!c.isSatisfied(day, materia, assigned, ctx)) return { passes: false, failedConstraint: c.name };
        }
        return { passes: true };
      }

      checkSlot(franja: FranjaHoraria, materia: MateriaInput, assigned: FranjaHoraria[], strict: boolean): boolean {
        return this.slotConstraints.every(c => c.isSatisfied(franja, materia, assigned, strict));
      }
    }
    \end{lstlisting}

    \subsection{Implementación: verificador (Student)}

    El verificador de horarios estudiantiles evalúa restricciones que operan sobre el conjunto de asignaciones confirmadas del estudiante:
    \begin{lstlisting}[language=TypeScript,caption={Verificador de restricciones (Student)}]
    export class ConstraintChecker {
      constructor(private readonly constraints: ScheduleConstraint[]) {}

      check(candidate: AsignacionResultado, confirmed: AsignacionResultado[], grupoElegido: Map<string, string>) {
        for (const constraint of this.constraints) {
          if (!constraint.isSatisfied(candidate, confirmed, grupoElegido)) {
            return { passes: false, failedConstraint: constraint.name };
          }
        }
        return { passes: true };
      }
    }
    \end{lstlisting}

    \subsection{Ejemplos concretos}

    Un par de implementaciones concretas, extraídas del código base, para ilustrar el patrón:

    \begin{lstlisting}[language=TypeScript,caption={Restricción concreta: preferencia de laboratorio por la tarde}]
    export class LabPrefersAfternoonConstraint extends SlotConstraint {
      name = "LabPrefersAfternoon [R7]";

      isSatisfied(franja: FranjaHoraria, materia: MateriaInput, _assigned: FranjaHoraria[], strict: boolean): boolean {
        if (!strict || materia.tipo !== "laboratorio") return true;
        const hora = parseInt(franja.inicio.split(":")[0] ?? "0", 10);
        return hora >= 10;
      }
    }
    \end{lstlisting}

    \begin{lstlisting}[language=TypeScript,caption={Restricción concreta: no solapamiento (student)}]
    export class NoOverlapConstraint extends ScheduleConstraint {
      name = "NoOverlap";

      isSatisfied(candidate: AsignacionResultado, confirmed: AsignacionResultado[]): boolean {
        return !confirmed.some(c => this.sesolapan(c, candidate));
      }
    }
    \end{lstlisting}

    \subsection{Buenas prácticas y organización}

    Para que el informe sea legible y útil en su conjunto sugiero:
    \begin{itemize}
      \item Mantener la distinción entre "allocator" y "student" en todos los ejemplos y diagramas.
      \item Documentar cada restricción con su propósito, ejemplos de entrada/salida y código mínimo (como se hizo arriba).
      \item Centralizar los diagramas en la carpeta \texttt{images/} y preferir rutas relativas para reproducibilidad; opcionalmente indicar URL alternativas.
    \end{itemize}

    Si quieres, puedo:
    \begin{enumerate}
      \item Cambiar todas las referencias de imágenes a URLs remotas.
      \item Extraer los listados de restricciones a anexos individuales (uno por cada restricción) para mantener la sección principal más concisa.
      \item Generar un índice detallado de clases/archivos del repositorio para que el lector pueda localizar la implementación exacta (por ejemplo: \texttt{src/universidad-proyecto/logic/schedule/...}).
    \end{enumerate}

    Dime cuál de estas opciones prefieres y la aplico.
  }

  checkSlot(franja: FranjaHoraria, materia: MateriaInput, assigned: FranjaHoraria[], strict: boolean): boolean {
    return this.slotConstraints.every(c => c.isSatisfied(franja, materia, assigned, strict));
  }
}
\end{lstlisting}

\subsection{Ejemplo concreto: preferencia de laboratorio por la tarde}
\begin{lstlisting}[language=TypeScript,caption={Restricción concreta de laboratorio}]
export class LabPrefersAfternoonConstraint extends SlotConstraint {
  name = "LabPrefersAfternoon [R7]";

  isSatisfied(franja: FranjaHoraria, materia: MateriaInput, _assigned: FranjaHoraria[], strict: boolean): boolean {
    if (!strict || materia.tipo !== "laboratorio") return true;
    const hora = parseInt(franja.inicio.split(":")[0] ?? "0", 10);
    return hora >= 10;
  }
}
\end{lstlisting}

\section{Constraints para horarios estudiantiles}

Los horarios de estudiante usan un verificador separado. La estructura es equivalente: una clase abstracta, clases concretas y un chequeo central.

\begin{lstlisting}[language=TypeScript,caption={Verificador de restricciones estudiantiles}]
export abstract class ScheduleConstraint {
  abstract name: string;
  abstract isSatisfied(candidate: AsignacionResultado, confirmed: AsignacionResultado[], grupoElegido: Map<string, string>): boolean;
}

export class ConstraintChecker {
  constructor(private readonly constraints: ScheduleConstraint[]) {}

  check(candidate: AsignacionResultado, confirmed: AsignacionResultado[], grupoElegido: Map<string, string>) {
    for (const constraint of this.constraints) {
      if (!constraint.isSatisfied(candidate, confirmed, grupoElegido)) {
        return { passes: false, failedConstraint: constraint.name };
      }
    }
    return { passes: true };
  }
}
\end{lstlisting}

\subsection{Ejemplo concreto: no solapamiento}
\begin{lstlisting}[language=TypeScript,caption={Restricción concreta de no solapamiento}]
export class NoOverlapConstraint extends ScheduleConstraint {
  name = "NoOverlap";

  isSatisfied(candidate: AsignacionResultado, confirmed: AsignacionResultado[]): boolean {
    return !confirmed.some(c => this.sesolapan(c, candidate));
  }
}
\end{lstlisting}

Con esta estructura, una nueva restricción se agrega creando una clase que implemente el contrato y registrándola en el verificador correspondiente.
